\section{The Architect's Question}\label{the-architects-question}

After this chapter we should be able to look at a characterized workload
and place it on the right \textbf{archetype} --- the canonical
deployment shapes that recur across the industry --- without reinventing
a design from scratch. We will map the four tiers (local/edge, server,
cluster, hyperscale), state what each can serve in model size and
traffic, and develop the decision rule for when a workload outgrows a
tier and must escalate. After this chapter, ``we run on a server'' or
``we need a cluster'' is a defensible, quantified statement tied to the
workload's numbers --- not a guess.

\section{Concept}\label{concept}

Architectures recur not because engineers lack imagination but because
the \emph{constraints} recur. Four archetypes cover the vast majority of
AI deployments, distinguished by how many GPUs are involved and over
what interconnect:

\begin{enumerate}
\def\labelenumi{\arabic{enumi}.}
\item
  \textbf{Local / edge} --- a single GPU (or a laptop-class
  accelerator). Serves small models (up to \textasciitilde7-13B at tight
  context) that must run on-device for privacy, latency, or connectivity
  reasons. Interconnect is not the question; memory capacity is.
\item
  \textbf{Server} --- one host with 1-8 GPUs, typically H100/H200,
  joined by NVLink/NVSwitch. Serves dense models up to
  \textasciitilde70B-180B with real concurrency. This is the canonical
  enterprise-Q\&A tier and the book's 8×H100 baseline. {[}2° FACT{]}
\item
  \textbf{Cluster} --- multiple hosts (tens to hundreds of GPUs) joined
  by InfiniBand or high-speed Ethernet. Required when a model exceeds a
  node's memory (so it must be split per Ch10) or when traffic outgrows
  one host (so P/D disaggregation per Ch11 applies).
\item
  \textbf{Hyperscale / data-center fleet} --- dedicated fleets of
  specialized inference/training clusters, often MoE or KV-centric
  (Mooncake-style disaggregation {[}S4{]}{[}1P{]}), serving many tenants
  at enormous scale. The economics favor extreme specialization.
\end{enumerate}

These are a \emph{continuum}, not hard categories; the useful question
is always \textbf{which tier does OUR workload's bounds demand?}

\section{Mental Model}\label{mental-model}

Think of each tier as a \textbf{box with a fixed memory budget and a
fixed cross-node bandwidth}. A workload fits a tier if (a) the model's
weight + KV residency fits the tier's total GPU memory (Ch7), and (b)
the tier's interconnect can carry the communication the chosen
parallelism requires (Ch9/Ch10), and (c) the resulting throughput +
latency meets the SLO (Ch11).

The mental model is a \textbf{ladder with explicit rungs}: we climb a
rung when a binding constraint can no longer be met. The three classic
triggers for climbing:

\begin{itemize}
\tightlist
\item
  \textbf{Memory}: model + KV no longer fits one node → cluster (split
  weights).
\item
  \textbf{Traffic}: one host's throughput ceiling (goodput) below demand
  → more replicas (DP across nodes) or P/D disaggregation → cluster.
\item
  \textbf{Latency/geography}: SLO can't be met due to distance or
  contention → move the serving edge closer (local) or split
  prefill/decode.
\end{itemize}

We do not climb for fashion; we climb when a measured number crosses a
threshold.

\section{Worked Example}\label{worked-example}

\subsection{Placing the Canonical
Workload}\label{placing-the-canonical-workload}

The canonical enterprise-Q\&A RAG workload sits firmly on the
\textbf{server} tier. We verify each bound:

\begin{itemize}
\item
  \textbf{Memory}: 70B FP16 weights = 140 GB; KV for 9.2K context at
  FP16 ≈ 24.7 GB; inference residency ≈ 164.7 GB {[}2° DERIVED{]}. An
  8×H100 host (640 GB) fits this with ample headroom for concurrency
  (Ch11 showed \textasciitilde24 GB per request × concurrent requests
  still fits). ✓ {[}2° DERIVED{]}
\item
  \textbf{Throughput}: \textasciitilde10 rps average / 40 rps peak,
  input-heavy (9.2K in / 300 out). Ch11 showed continuous batching +
  prefix caching keeps p95 TTFT under the 2 s SLO on a single 8×H100
  host. ✓
\item
  \textbf{Interconnect}: parallelism is unnecessary (one node) --- the
  interconnect question doesn't bind. ✓
\end{itemize}

So a single 8×H100 server is the right archetype --- which is exactly
what earlier chapters assumed. {[}2° DERIVED{]}

\subsection{Escalating: When 10× Traffic Forces the Cluster
Tier}\label{escalating-when-10-traffic-forces-the-cluster-tier}

Now suppose the workload's demand grows to \textasciitilde100 rps
average (a 10×). Recompute the ceiling:

\begin{itemize}
\item
  \textbf{Throughput ceiling of one host}: continuous batching on 8×H100
  at \textasciitilde9.5K tokens/request sustained roughly tens of
  thousands of tokens/s; goodput against the SLO is the binding number.
  At \textasciitilde100 rps × 9.5K tokens = \textasciitilde950K
  tokens/s, one host's goodput is exceeded. {[}2° DERIVED{]}
\item
  \textbf{Response}: replicate the model across multiple servers (data
  parallelism at the serving layer --- run N identical 8×H100 hosts
  behind a load balancer), so each host carries \textasciitilde10 rps
  again. This is the \textbf{cluster} tier: the model still fits one
  node (so no weight-splitting), but the \emph{fleet} of replicas is a
  small cluster. The escalation is driven by traffic, not memory. {[}2°
  DERIVED{]}
\end{itemize}

\subsection{Dropping a Tier: When Privacy Forces
Local}\label{dropping-a-tier-when-privacy-forces-local}

If the workload must run fully on-site without data leaving a facility,
and the model can be distilled to a 7-13B capable of the Q\&A task, the
\textbf{local/edge} tier applies: a single GPU runs the distilled model
at reduced (but acceptable) accuracy for privacy. The tier choice traded
capability for the operational constraint (privacy). {[}2° DERIVED{]}

\subsubsection{\hyperref[tab:13.1]{Table 13-1} --- Reference architecture
archetypes}\label{table-13-1-reference-architecture-archetypes}

\begin{longtable}[]{@{}
  >{\raggedright\arraybackslash}p{0.1667\linewidth}
  >{\raggedright\arraybackslash}p{0.1667\linewidth}
  >{\raggedright\arraybackslash}p{0.1667\linewidth}
  >{\raggedright\arraybackslash}p{0.1667\linewidth}
  >{\raggedright\arraybackslash}p{0.1667\linewidth}
  >{\raggedright\arraybackslash}p{0.1667\linewidth}@{}}
\toprule\noalign{}
\begin{minipage}[b]{\linewidth}\raggedright
Tier
\end{minipage} & \begin{minipage}[b]{\linewidth}\raggedright
GPUs
\end{minipage} & \begin{minipage}[b]{\linewidth}\raggedright
Interconnect
\end{minipage} & \begin{minipage}[b]{\linewidth}\raggedright
Model range
\end{minipage} & \begin{minipage}[b]{\linewidth}\raggedright
Typical latency
\end{minipage} & \begin{minipage}[b]{\linewidth}\raggedright
When to use
\end{minipage} \\
\midrule\noalign{}
\endhead
\bottomrule\noalign{}
\endlastfoot
Local / edge & 1 & --- & ≤7-13B dense & low (on-device) & privacy,
offline, connectivity \\
Server & 1-8 & NVLink/NVSwitch & ≤\textasciitilde70-180B dense & meets
§14 SLO & most single-tenant enterprise \\
Cluster & tens-hundreds & InfiniBand/Eth & any that needs splitting &
depends & model \textgreater{} node, or traffic \textgreater{} host
goodput \\
Hyperscale fleet & dedicated & high-bandwidth fabric & MoE+, specialized
& extreme-scale economics & multi-tenant, largest scale \\

\label{tab:13.1}\end{longtable}

\emph{(Bounds are {[}2° FACT{]}/{[}2° DERIVED{]} design heuristics, not
vendor guarantees.)}

\section{Measurement}\label{measurement}

The tier decision is only as good as the numbers that trigger it. We
measure three things:

\begin{enumerate}
\def\labelenumi{\arabic{enumi}.}
\tightlist
\item
  \textbf{Residency vs tier memory} (Ch7): does weight + KV fit? If not,
  the floor is a cluster.
\item
  \textbf{Host goodput ceiling vs demand} (Ch6/Ch11): measure the max
  rps/tokens/s the host delivers at the SLO, and compare to demand. When
  demand \textgreater{} ceiling, escalate.
\item
  \textbf{Communication fraction} (Ch9/Ch10): if the parallelism the
  cluster requires spends \textgreater\textasciitilde20\% of step time
  on collectives, the interconnect (not the GPU) is the binding tier
  constraint.
\end{enumerate}

These three numbers locate the workload on the ladder objectively.

\section{Common Mistakes}\label{common-mistakes}

\begin{itemize}
\tightlist
\item
  \textbf{Choosing hyperscale before measuring goodput.} Most workloads
  never exceed one host's goodput; building a fleet first is expensive
  and unnecessary.
\item
  \textbf{Using model size as the only tier driver.} A 70B model fits a
  server; traffic, not size, usually forces the cluster --- an architect
  who only watches parameters escalates too early or too late.
\item
  \textbf{Ignoring the operational constraint.} Privacy/geography can
  force a \emph{lower} tier regardless of what compute wants; the ladder
  is bounded by operations, not just performance.
\item
  \textbf{Assuming a smaller model can't serve the tier.} Distilled
  small models on local GPUs can meet the goal when the capability
  suffices --- don't over-provision the tier.
\item
  \textbf{Treating tiers as walls.} They are a continuum; the correct
  answer can be a hybrid (edge cache + server generation) the measured
  numbers justify.
\end{itemize}

\section{Architecture Consequence}\label{architecture-consequence}

The archetype sets the entire downstream design vocabulary. On the
\textbf{server} tier (canonical), the consequences are already
established: continuous batching + PagedAttention + prefix caching
(Ch11), KV quantization when residency binds (Ch7), no parallelism
(Ch10). Moving up to \textbf{cluster} introduces: DP/PP/TP replicas
(Ch10), P/D disaggregation (Ch11), and interconnect budgeting (Ch9).
Moving \textbf{down} to local forces distillation and on-device KV
limits. The tier is the top-level decision; every other choice in the
book slots beneath it.

\begin{figure}
\centering
\pandocbounded{\includegraphics[keepaspectratio,alt={\hyperref[fig:13.1]{Fig 13.1} --- The reference-architecture ladder and its escalation triggers {[}ILLUSTRATIVE conceptual{]}},width=\maxwidth{\textwidth},keepaspectratio]{figures/fig-13-1301.pdf}}
\caption{\hyperref[fig:13.1]{Fig 13.1} --- The reference-architecture ladder and its
escalation triggers {[}ILLUSTRATIVE conceptual{]}}

\label{fig:13.1}\end{figure}

\emph{\hyperref[fig:13.1]{Fig 13.1} --- The reference-architecture ladder: single GPU →
single host → multi-host → cluster/fleet, with escalation triggers (red:
model/KV/throughput outgrows a tier) and de-escalation (green: privacy,
data sovereignty, cost) moving a workload between tiers.}

\section{What We Still Don't Know}\label{what-we-still-dont-know}

\begin{itemize}
\tightlist
\item
  \textbf{Exact host goodput ceilings} depend on the specific server,
  model, degree and batching; the representative numbers here are
  {[}VERIFY{]} per deployment.
\item
  \textbf{When P/D disaggregation's benefit crosses its cost} for
  mid-tier workloads is {[}HYPOTHESIS{]}; the breakpoint is workload-
  and fabric-specific.
\item
  \textbf{Fleet-level economics} (Ch16 TCO across a cluster) interact
  with the tier choice in ways that are {[}HYPOTHESIS{]} until a real
  cost model is applied.
\end{itemize}

\section{End-of-Chapter Mini-Case}\label{end-of-chapter-mini-case}

An architect returns to the canonical internal Q\&A service a year
later. The company has grown from \textasciitilde2,000 to
\textasciitilde20,000 users; traffic is now \textasciitilde10× (≈100 rps
average). The on-prem service is saturated: p95 TTFT has drifted to 3 s
against the 2 s SLO. Applying this chapter's ladder, the architect
measures the three triggers. Memory: the 70B model still fits one 8×H100
node (164.7 GB residency), so the \emph{model} does not force a cluster.
Traffic: host goodput is exceeded at \textasciitilde100 rps --- the
binding trigger. Operational: the data privacy requirement still forbids
a public hyperscale cloud. The architect's decision: escalate to the
\textbf{cluster} tier but keep it on-prem --- run a small fleet of
8×H100 replicas behind a load balancer so each carries \textasciitilde10
rps again, staying within both the memory and the privacy constraints.
The result: p95 TTFT back under 1.8 s at 100 rps, on site, without
hyperscale --- because the architect climbed the ladder precisely one
rung, and the correct rung was determined by the measured traffic
trigger, not by fashion or by model size alone.
